\documentclass[11pt,a4paper]{article}

\usepackage[T1]{fontenc}
\usepackage[utf8]{inputenc}
\usepackage{lmodern}
\usepackage[a4paper,hmargin=26mm,vmargin=24mm]{geometry}
\usepackage{amsmath,amssymb,amsthm,mathtools}
\usepackage{stmaryrd}
\usepackage{microtype}
\usepackage{enumitem}
\usepackage{booktabs}
\PassOptionsToPackage{hyphens}{url}
\usepackage[hidelinks]{hyperref}
\usepackage[capitalise,noabbrev]{cleveref}
\crefname{equation}{}{}
\Crefname{equation}{Equation}{Equations}
\creflabelformat{equation}{(#2#1#3)}

\theoremstyle{plain}
\newtheorem{theorem}{Theorem}
\newtheorem{lemma}[theorem]{Lemma}
\newtheorem{proposition}[theorem]{Proposition}
\newtheorem{corollary}[theorem]{Corollary}
\theoremstyle{definition}
\newtheorem{definition}[theorem]{Definition}
\theoremstyle{remark}
\newtheorem{remark}[theorem]{Remark}

\DeclareMathOperator{\Fb}{Fb}
\DeclareMathOperator{\Tr}{Tr}
\DeclareMathOperator{\Cut}{Cut}
\newcommand{\R}{\mathbb{R}}
\newcommand{\Rs}{\mathbb{R}(s)}
\newcommand{\id}{\mathrm{id}}
\newcommand{\then}{\mathbin{\fatsemi}}
\newcommand{\Mat}{\mathbf{Mat}}
\newcommand{\LTIq}{\mathbf{LTI}/{\sim}}
\newcommand{\MealyMod}{\mathbf{Mealy}/{\sim}}
\newcommand{\sys}[5]{\bigl\langle #1,\ #2,\ #3,\ #4,\ #5 \bigr\rangle}
\newcommand{\lbl}[1]{\nolinkurl{#1}}
\setcounter{MaxMatrixCols}{12}

\title{Feedback operators on matrices:\\ the cut, guarded partial traces, and the transfer function}
\author{Working note for the ACT4E chapter on feedback}
\date{24 November 2026}

\begin{document}
\maketitle

\noindent
This note is a sequel to the note \emph{The feedback operator on $\LTIq$ is a partial trace guarded by the delay} (file \lbl{LTI-feedback-vs-partial-trace}), which we call the first note.
It answers two questions about the category $\Mat_\R$ of real matrices, with block-diagonal matrices as monoidal product (\lbl{exa:MatR-strict-monoidal}), and about its partial trace (\lbl{def:trace-for-mat}):
\begin{quote}
    Is there a feedback operator on $\Mat_\R$ which is analogous to the feedback operators of $\MealyMod$ and of $\LTIq$?
    Can it be related to the partial trace through a delay?
\end{quote}

The feedback operators of $\MealyMod$ and of $\LTIq$ move a wire into the state.
A matrix has no state, so this construction cannot be carried out inside $\Mat_\R$.
The correct analogy is a different one: real matrices are to $\LTIq$ what functions are to $\MealyMod$.
A real matrix is the same thing as an LTI system with state space $\R^0$, and applying the feedback operator $\Fb^p$ of the first note to such a system gives a system with state space $\R^p$, which is not of this kind if $p \geq 1$.

One can nevertheless ask which feedback operators $\Mat_\R$ has.
The results are as follows.
\begin{enumerate}[nosep]
    \item ``Cutting the loop'' is a feedback operator on $\Mat_\R$ (\cref{prop:cut}).
    \item It is the only one (\cref{thm:unique}).
          In particular, $\Mat_\R$ has no trace operator (\cref{cor:no-trace}).
          The delay of the cut is zero, and the cut is the partial trace guarded by this delay (\cref{cor:cut-delay}).
    \item For every scalar $\lambda$, the partial trace guarded by $\lambda$ is a partially defined operator $\Fb_\lambda$ whose delay is $\lambda$ times the identity (\cref{lem:guard,lem:delay}).
          The partial trace is $\Fb_1$ and the cut is $\Fb_0$.
          The argument for uniqueness shows that there is no room for other operators of this kind, with one exception (\cref{rem:classification}).
    \item The transfer function of an LTI system is $\Fb_{1/s}$ applied to the matrix of the system (\cref{prop:evaluation}), and it turns the feedback operator of $\LTIq$ into $\Fb_{1/s}$ (\cref{thm:transfer}).
\end{enumerate}
What this note does \emph{not} prove is listed in \cref{rem:not-proved}.

\section{Setup}

\paragraph{Scalars and matrices.}
Throughout, $K$ is a field.
The case of the book is $K = \R$; in \cref{sec:transfer} we also need the field $K = \Rs$ of rational functions in one indeterminate $s$ with real coefficients.
The letters $n$, $m$, $p$, $q$, $e$, $k$ denote natural numbers, and each of them may be zero.
We write $K^{m \times n}$ for the set of matrices with $m$ rows, $n$ columns, and entries in $K$.
We allow matrices with zero rows or with zero columns.
We write $I_p$ for the identity matrix of size $p \times p$, and $0$ for a zero matrix whose size is determined by the context.
In a block matrix, a block with zero rows or zero columns is simply absent.
The matrix $I_0$, which has no entries, is invertible.
We use the following fact: a square matrix with entries in a field is invertible if and only if the zero vector is the only vector which it maps to zero.

\paragraph{The category of matrices.}
The category $\Mat_K$ is defined as in the book for $K = \R$ (\lbl{def:MatR}, \lbl{exa:MatR-strict-monoidal}).
Its objects are the coordinate spaces $K^n$.
A morphism $M \colon K^n \to K^m$ is a matrix $M \in K^{m \times n}$.
Composition is the matrix product in the opposite order, the identity of $K^n$ is $I_n$, and the monoidal product is the block-diagonal matrix:
\begin{equation}
    \label{eq:mat-operations}
    M \then N = N M,
    \qquad
    \id_{K^n} = I_n,
    \qquad
    M \oplus N = \begin{bmatrix} M & 0 \\ 0 & N \end{bmatrix} .
\end{equation}
The isomorphisms of $\Mat_K$ are the invertible matrices; they are square.
The braiding of $K^p$ with itself is the matrix (compare \lbl{eq:LTI-swap-matrix})
\begin{equation}
    \label{eq:swap}
    S_{p,p} = \begin{bmatrix} 0 & I_p \\ I_p & 0 \end{bmatrix} \colon K^{p+p} \to K^{p+p} .
\end{equation}
The braiding does not occur in the axioms of a feedback operator.
We need it only to speak about delays.

\paragraph{Blocks.}
Let $M \in K^{(m+p) \times (n+p)}$.
The \emph{blocks of $M$ with respect to $p$} are the matrices in the decomposition
\begin{equation}
    \label{eq:blocks}
    M = \begin{bmatrix} M_{11} & M_{12} \\ M_{21} & M_{22} \end{bmatrix},
    \qquad
    M_{11} \in K^{m \times n}, \quad
    M_{12} \in K^{m \times p}, \quad
    M_{21} \in K^{p \times n}, \quad
    M_{22} \in K^{p \times p} .
\end{equation}
For $p = 1$ we write the decomposition as
\begin{equation}
    \label{eq:blocks-one}
    M = \begin{bmatrix} M_{11} & c \\ r & d \end{bmatrix},
\end{equation}
where $c \in K^{m \times 1}$ is a column, $r \in K^{1 \times n}$ is a row, and $d \in K$.

\begin{definition}[Feedback operator on $\Mat_K$]
    \label{def:fb-mat}
    A \emph{feedback operator} on $\Mat_K$ is a family of functions
    \begin{equation}
        \Fb^p_{n,m} \colon K^{(m+p) \times (n+p)} \to K^{m \times n},
    \end{equation}
    one for every triple $n$, $m$, $p$, which satisfies the following conditions.
    We write $\Fb^p$ for $\Fb^p_{n,m}$ when $n$ and $m$ are clear from the matrix to which it is applied.
    \begin{enumerate}[label=\textup{(F\arabic*)}, leftmargin=*]
        \item \emph{Naturality in the source.}
              For all $M \in K^{(m+p) \times (n+p)}$ and $G \in K^{n \times n'}$,
              \begin{equation*}
                  \Fb^p\bigl( M (G \oplus I_p) \bigr) = \Fb^p(M)\, G .
              \end{equation*}
        \item \emph{Naturality in the target.}
              For all $M \in K^{(m+p) \times (n+p)}$ and $H \in K^{m' \times m}$,
              \begin{equation*}
                  \Fb^p\bigl( (H \oplus I_p) M \bigr) = H\, \Fb^p(M) .
              \end{equation*}
        \item \emph{Dinaturality with respect to isomorphisms.}
              For all $M \in K^{(m+p) \times (n+p)}$ and all invertible $h \in K^{p \times p}$,
              \begin{equation*}
                  \Fb^p\bigl( (I_m \oplus h) M \bigr) = \Fb^p\bigl( M (I_n \oplus h) \bigr) .
              \end{equation*}
        \item \emph{Vanishing I.}
              For all $M \in K^{m \times n}$, we have $\Fb^0(M) = M$.
        \item \emph{Vanishing II.}
              For all $M \in K^{(m+p+q) \times (n+p+q)}$,
              \begin{equation*}
                  \Fb^p_{n,m}\bigl( \Fb^q_{n+p,\,m+p}(M) \bigr) = \Fb^{p+q}_{n,m}(M) .
              \end{equation*}
        \item \emph{Superposing.}
              For all $M \in K^{(m+p) \times (n+p)}$ and all $e$,
              \begin{equation*}
                  \Fb^p_{e+n,\,e+m}( I_e \oplus M ) = I_e \oplus \Fb^p_{n,m}(M) .
              \end{equation*}
    \end{enumerate}
\end{definition}

\begin{remark}
    \label{rem:axioms}
    The conditions (F1) to (F6) are the six axioms of \lbl{def:feedback-cat}, in the same order, written out for $\Mat_K$ by means of \cref{eq:mat-operations}.
    For instance, the first axiom reads $\Fb\bigl( (g \otimes \id) \then f \bigr) = g \then \Fb(f)$; for $f = M$ and $g = G$ we have $(g \otimes \id) \then f = M (G \oplus I_p)$ and $g \then \Fb(f) = \Fb(M)\, G$.
    In the third axiom, the isomorphism is a morphism $K^{p'} \to K^{p}$; since invertible matrices are square, $p' = p$.
\end{remark}

\begin{definition}[Partial trace]
    \label{def:tr}
    A matrix $M \in K^{(m+p) \times (n+p)}$, with blocks as in \cref{eq:blocks}, is \emph{$p$-traceable} if $I_p - M_{22}$ is invertible.
    In this case,
    \begin{equation}
        \label{eq:tr}
        \Tr^p(M) \coloneqq M_{11} + M_{12} (I_p - M_{22})^{-1} M_{21} \in K^{m \times n} .
    \end{equation}
    For $K = \R$ this is \lbl{def:trace-for-mat}, where the blocks $M_{11}$, $M_{12}$, $M_{21}$, $M_{22}$ are called $A$, $C$, $B$, $D$.
\end{definition}

\section{Cutting the loop is a feedback operator}

\begin{definition}[Cut]
    \label{def:cut}
    For $M \in K^{(m+p) \times (n+p)}$, with blocks as in \cref{eq:blocks}, we set
    \begin{equation}
        \Cut^p(M) \coloneqq M_{11} \in K^{m \times n} .
    \end{equation}
    So $\Cut^p(M)$ consists of the first $m$ rows and the first $n$ columns of $M$.
\end{definition}

\begin{proposition}
    \label{prop:cut}
    The family $\Cut^p$ is a feedback operator on $\Mat_K$.
\end{proposition}

\begin{proof}
    Let $M \in K^{(m+p) \times (n+p)}$, with blocks as in \cref{eq:blocks}.

    (F1) and (F2).
    For $G \in K^{n \times n'}$ and $H \in K^{m' \times m}$,
    \begin{equation*}
        M (G \oplus I_p) = \begin{bmatrix} M_{11} G & M_{12} \\ M_{21} G & M_{22} \end{bmatrix},
        \qquad
        (H \oplus I_p) M = \begin{bmatrix} H M_{11} & H M_{12} \\ M_{21} & M_{22} \end{bmatrix},
    \end{equation*}
    and these are the decompositions with respect to $p$.
    So $\Cut^p$ of the first matrix is $M_{11} G = \Cut^p(M)\, G$, and $\Cut^p$ of the second one is $H M_{11} = H \Cut^p(M)$.

    (F3).
    For $h \in K^{p \times p}$,
    \begin{equation*}
        (I_m \oplus h) M = \begin{bmatrix} M_{11} & M_{12} \\ h M_{21} & h M_{22} \end{bmatrix},
        \qquad
        M (I_n \oplus h) = \begin{bmatrix} M_{11} & M_{12} h \\ M_{21} & M_{22} h \end{bmatrix} .
    \end{equation*}
    Both matrices have $M_{11}$ as their block in the upper left, so $\Cut^p$ takes the same value on both.
    Here we did not use that $h$ is invertible.

    (F4).
    For $p = 0$, the blocks $M_{12}$, $M_{21}$, $M_{22}$ are absent, and $M = M_{11} = \Cut^0(M)$.

    (F5).
    Let $M \in K^{(m+p+q) \times (n+p+q)}$.
    The matrix $\Cut^q(M)$ consists of the first $m+p$ rows and the first $n+p$ columns of $M$, and $\Cut^p(\Cut^q(M))$ consists of the first $m$ rows and the first $n$ columns of that matrix.
    These are the first $m$ rows and the first $n$ columns of $M$, that is, $\Cut^{p+q}(M)$.

    (F6).
    The decomposition of $I_e \oplus M$ with respect to $p$ is
    \begin{equation*}
        I_e \oplus M =
        \left[\begin{array}{cc|c}
                I_e & 0      & 0      \\
                0   & M_{11} & M_{12} \\ \hline
                0   & M_{21} & M_{22}
            \end{array}\right],
    \end{equation*}
    so $\Cut^p(I_e \oplus M) = \begin{bsmallmatrix} I_e & 0 \\ 0 & M_{11} \end{bsmallmatrix} = I_e \oplus \Cut^p(M)$.
\end{proof}

\section{The cut is the only feedback operator}
\label{sec:unique}

In this section, $\Fb$ denotes an arbitrary feedback operator on $\Mat_K$.
The proof that $\Fb = \Cut$ has three steps.
First, for one wire, $\Fb$ is determined by a single function $\varphi \colon K \to K$ (\cref{lem:one-wire}).
Second, comparing two wires with the same two wires after a change of coordinates gives an equation for $\varphi$ (\cref{lem:functional-equation}).
Third, this equation has no solution which is defined on all of $K$, except $\varphi = 0$ (\cref{lem:solutions}).

For $d \in K$ we consider the matrix
\begin{equation}
    \label{eq:Kd}
    K_d \coloneqq \begin{bmatrix} 0 & 1 \\ 1 & d \end{bmatrix} \in K^{(1+1) \times (1+1)} .
\end{equation}
The matrix $\Fb^1_{1,1}(K_d)$ has one row and one column.
We define $\varphi(d) \in K$ to be its entry:
\begin{equation}
    \label{eq:phi}
    \Fb^1_{1,1}(K_d) = \begin{bmatrix} \varphi(d) \end{bmatrix} .
\end{equation}

\begin{lemma}[One wire]
    \label{lem:one-wire}
    For all $n$ and $m$, and for every $M \in K^{(m+1) \times (n+1)}$ with blocks as in \cref{eq:blocks-one},
    \begin{equation}
        \label{eq:one-wire}
        \Fb^1(M) = M_{11} + \varphi(d)\, c\, r .
    \end{equation}
\end{lemma}

\begin{proof}
    Consider the matrices
    \begin{equation*}
        G \coloneqq \begin{bmatrix} I_n \\ r \end{bmatrix} \in K^{(n+1) \times n},
        \qquad
        H \coloneqq \begin{bmatrix} M_{11} & c \end{bmatrix} \in K^{m \times (n+1)} .
    \end{equation*}
    We claim that
    \begin{equation}
        \label{eq:factorization}
        M = (H \oplus I_1)\, (I_n \oplus K_d)\, (G \oplus I_1) .
    \end{equation}
    Indeed, the three factors are
    \begin{equation*}
        H \oplus I_1 = \begin{bmatrix} M_{11} & c & 0 \\ 0 & 0 & 1 \end{bmatrix},
        \qquad
        I_n \oplus K_d = \begin{bmatrix} I_n & 0 & 0 \\ 0 & 0 & 1 \\ 0 & 1 & d \end{bmatrix},
        \qquad
        G \oplus I_1 = \begin{bmatrix} I_n & 0 \\ r & 0 \\ 0 & 1 \end{bmatrix},
    \end{equation*}
    where the columns of the first two matrices and the rows of the last two are grouped as $n + 1 + 1$.
    We compute
    \begin{equation*}
        (I_n \oplus K_d)(G \oplus I_1) = \begin{bmatrix} I_n & 0 \\ 0 & 1 \\ r & d \end{bmatrix},
        \qquad
        (H \oplus I_1) \begin{bmatrix} I_n & 0 \\ 0 & 1 \\ r & d \end{bmatrix} = \begin{bmatrix} M_{11} & c \\ r & d \end{bmatrix} = M .
    \end{equation*}
    This proves \cref{eq:factorization}.
    Now we apply $\Fb^1$:
    \begin{align*}
        \Fb^1(M)
         & = \Fb^1\bigl( (H \oplus I_1)\, (I_n \oplus K_d)\, (G \oplus I_1) \bigr) \\
         & = H\, \Fb^1\bigl( (I_n \oplus K_d)\, (G \oplus I_1) \bigr)               &  & \text{by (F2)} \\
         & = H\, \Fb^1( I_n \oplus K_d )\, G                                        &  & \text{by (F1)} \\
         & = H\, \bigl( I_n \oplus \Fb^1(K_d) \bigr)\, G                            &  & \text{by (F6)} \\
         & = \begin{bmatrix} M_{11} & c \end{bmatrix}
        \begin{bmatrix} I_n & 0 \\ 0 & \varphi(d) \end{bmatrix}
        \begin{bmatrix} I_n \\ r \end{bmatrix}                                     &  & \text{by \cref{eq:phi}} \\
         & = M_{11} + \varphi(d)\, c\, r .
    \end{align*}
    In the second line, (F2) is applied to the matrix $(I_n \oplus K_d)(G \oplus I_1) \in K^{((n+1)+1) \times (n+1)}$.
    In the third line, (F1) is applied to the matrix $I_n \oplus K_d \in K^{((n+1)+1) \times ((n+1)+1)}$.
    In the fourth line, (F6) is applied to $K_d$, with $e = n$.
\end{proof}

\begin{lemma}[The functional equation]
    \label{lem:functional-equation}
    For all $a, b \in K$,
    \begin{equation}
        \label{eq:functional-equation}
        \varphi(a)\, \bigl( 1 + (b - a)\, \varphi(b) \bigr) = \varphi(b) .
    \end{equation}
\end{lemma}

\begin{proof}
    Let $a, b \in K$.
    Consider the matrices
    \begin{equation*}
        N \coloneqq \begin{bmatrix} 0 & 1 & 0 \\ 0 & a & -a \\ 1 & 0 & b \end{bmatrix} \in K^{(1+2) \times (1+2)},
        \qquad
        h \coloneqq \begin{bmatrix} 1 & 1 \\ 0 & 1 \end{bmatrix} \in K^{2 \times 2} .
    \end{equation*}
    The matrix $h$ is invertible.
    Multiplying $N$ from the left by $I_1 \oplus h$ adds the third row to the second row, and multiplying $N$ from the right by $I_1 \oplus h$ adds the second column to the third column:
    \begin{equation*}
        L \coloneqq (I_1 \oplus h) N = \begin{bmatrix} 0 & 1 & 0 \\ 1 & a & b - a \\ 1 & 0 & b \end{bmatrix},
        \qquad
        R \coloneqq N (I_1 \oplus h) = \begin{bmatrix} 0 & 1 & 1 \\ 0 & a & 0 \\ 1 & 0 & b \end{bmatrix} .
    \end{equation*}
    By (F3), with $n = m = 1$ and $p = 2$,
    \begin{equation}
        \label{eq:L-equals-R}
        \Fb^2(L) = \Fb^2(R) .
    \end{equation}
    We compute both sides.
    By (F5), with $n = m = p = q = 1$, for every $X \in K^{3 \times 3}$ we have $\Fb^2_{1,1}(X) = \Fb^1_{1,1}\bigl( \Fb^1_{2,2}(X) \bigr)$, and both applications of $\Fb^1$ are given by \cref{lem:one-wire}.

    \emph{The matrix $L$.}
    Its decomposition \cref{eq:blocks-one}, as an element of $K^{(2+1) \times (2+1)}$, has
    \begin{equation*}
        M_{11} = \begin{bmatrix} 0 & 1 \\ 1 & a \end{bmatrix},
        \qquad
        c = \begin{bmatrix} 0 \\ b - a \end{bmatrix},
        \qquad
        r = \begin{bmatrix} 1 & 0 \end{bmatrix},
        \qquad
        d = b,
        \qquad \text{so} \qquad
        c\, r = \begin{bmatrix} 0 & 0 \\ b - a & 0 \end{bmatrix} .
    \end{equation*}
    By \cref{eq:one-wire},
    \begin{equation*}
        \Fb^1_{2,2}(L) = \begin{bmatrix} 0 & 1 \\ 1 & a \end{bmatrix} + \varphi(b) \begin{bmatrix} 0 & 0 \\ b - a & 0 \end{bmatrix}
        = \begin{bmatrix} 0 & 1 \\ 1 + (b - a) \varphi(b) & a \end{bmatrix} .
    \end{equation*}
    The decomposition \cref{eq:blocks-one} of this matrix, as an element of $K^{(1+1) \times (1+1)}$, has $M_{11} = 0$, $c = 1$, $r = 1 + (b - a) \varphi(b)$, and $d = a$.
    By \cref{eq:one-wire},
    \begin{equation}
        \label{eq:Fb-L}
        \Fb^2(L) = 0 + \varphi(a) \cdot 1 \cdot \bigl( 1 + (b - a) \varphi(b) \bigr) = \varphi(a)\, \bigl( 1 + (b - a)\, \varphi(b) \bigr) .
    \end{equation}

    \emph{The matrix $R$.}
    Its decomposition \cref{eq:blocks-one}, as an element of $K^{(2+1) \times (2+1)}$, has
    \begin{equation*}
        M_{11} = \begin{bmatrix} 0 & 1 \\ 0 & a \end{bmatrix},
        \qquad
        c = \begin{bmatrix} 1 \\ 0 \end{bmatrix},
        \qquad
        r = \begin{bmatrix} 1 & 0 \end{bmatrix},
        \qquad
        d = b,
        \qquad \text{so} \qquad
        c\, r = \begin{bmatrix} 1 & 0 \\ 0 & 0 \end{bmatrix} .
    \end{equation*}
    By \cref{eq:one-wire},
    \begin{equation*}
        \Fb^1_{2,2}(R) = \begin{bmatrix} 0 & 1 \\ 0 & a \end{bmatrix} + \varphi(b) \begin{bmatrix} 1 & 0 \\ 0 & 0 \end{bmatrix}
        = \begin{bmatrix} \varphi(b) & 1 \\ 0 & a \end{bmatrix} .
    \end{equation*}
    The decomposition \cref{eq:blocks-one} of this matrix has $M_{11} = \varphi(b)$, $c = 1$, $r = 0$, and $d = a$.
    By \cref{eq:one-wire},
    \begin{equation}
        \label{eq:Fb-R}
        \Fb^2(R) = \varphi(b) + \varphi(a) \cdot 1 \cdot 0 = \varphi(b) .
    \end{equation}
    Here we identify a matrix with one row and one column with its entry.
    Inserting \cref{eq:Fb-L} and \cref{eq:Fb-R} into \cref{eq:L-equals-R} gives \cref{eq:functional-equation}.
\end{proof}

\begin{remark}
    \label{rem:conjugation}
    The matrices $L$ and $R$ were found as follows.
    The loop block of $R$ is the diagonal matrix with entries $a$ and $b$: in $R$, the two wires of the loop do not interact.
    The loop block of $L$ is $h \begin{bsmallmatrix} a & 0 \\ 0 & b \end{bsmallmatrix} h^{-1}$: this is the same loop, described in other coordinates.
    Condition (F3) says that the result of feeding back may not depend on the coordinates of the loop.
\end{remark}

The next lemma solves \cref{eq:functional-equation}.
We state it for functions which are defined only on a subset of $K$, because we need this in \cref{sec:guarded}.

\begin{lemma}[Solutions of the functional equation]
    \label{lem:solutions}
    Let $U \subseteq K$ be a subset, and let $\psi \colon U \to K$ be a function.
    The following are equivalent.
    \begin{enumerate}[label=\textup{(\roman*)}]
        \item For all $a, b \in U$, we have $\psi(a)\, \bigl( 1 + (b - a)\, \psi(b) \bigr) = \psi(b)$.
        \item Either $\psi(a) = 0$ for all $a \in U$, or there is an element $c \in K$ with $c \notin U$ such that
              \begin{equation*}
                  \psi(a) = \frac{1}{c - a} \qquad \text{for all } a \in U .
              \end{equation*}
    \end{enumerate}
    In particular, if $U = K$, then \textup{(i)} holds if and only if $\psi = 0$.
\end{lemma}

\begin{proof}
    (i) $\Rightarrow$ (ii).
    Suppose first that $\psi(a_0) = 0$ for some $a_0 \in U$.
    Then (i), with $a = a_0$, gives $0 = \psi(b)$ for all $b \in U$.
    So $\psi$ is zero.

    Suppose now that $\psi(a) \neq 0$ for all $a \in U$.
    If $U$ is empty, then $\psi$ is zero and there is nothing to prove.
    Otherwise choose $a_0 \in U$ and set $c \coloneqq a_0 + \psi(a_0)^{-1}$.
    Let $b \in U$.
    By (i), with $a = a_0$,
    \begin{equation*}
        \psi(a_0) = \psi(b) - (b - a_0)\, \psi(a_0)\, \psi(b) = \psi(b)\, \bigl( 1 - (b - a_0)\, \psi(a_0) \bigr) .
    \end{equation*}
    Dividing by $\psi(a_0)$ gives
    \begin{equation*}
        1 = \psi(b)\, \bigl( \psi(a_0)^{-1} - b + a_0 \bigr) = \psi(b)\, (c - b) .
    \end{equation*}
    So $c - b \neq 0$ and $\psi(b) = \frac{1}{c - b}$.
    This holds for every $b \in U$; in particular $c \neq b$ for every $b \in U$, that is, $c \notin U$.

    (ii) $\Rightarrow$ (i).
    If $\psi$ is zero, both sides of the equation in (i) are zero.
    If $\psi(a) = \frac{1}{c - a}$ for all $a \in U$, then for $a, b \in U$
    \begin{equation*}
        \psi(a)\, \bigl( 1 + (b - a)\, \psi(b) \bigr)
        = \frac{1}{c - a} \cdot \frac{(c - b) + (b - a)}{c - b}
        = \frac{1}{c - a} \cdot \frac{c - a}{c - b}
        = \psi(b) .
    \end{equation*}

    The last statement: if $U = K$, there is no element $c \in K$ with $c \notin U$, so the second case of (ii) cannot occur.
\end{proof}

\begin{theorem}[Uniqueness]
    \label{thm:unique}
    Let $K$ be a field, and let $\Fb$ be a feedback operator on $\Mat_K$.
    Then $\Fb = \Cut$: for all $n$, $m$, $p$ and all $M \in K^{(m+p) \times (n+p)}$, with blocks as in \cref{eq:blocks},
    \begin{equation}
        \Fb^p(M) = M_{11} .
    \end{equation}
\end{theorem}

\begin{proof}
    Let $\varphi \colon K \to K$ be the function defined by \cref{eq:phi}.
    By \cref{lem:functional-equation}, it satisfies condition (i) of \cref{lem:solutions}, with $U = K$.
    By the last statement of \cref{lem:solutions}, $\varphi = 0$.
    So \cref{lem:one-wire} gives, for all $n$ and $m$,
    \begin{equation}
        \label{eq:one-wire-cut}
        \Fb^1(M) = M_{11} = \Cut^1(M) \qquad \text{for all } M \in K^{(m+1) \times (n+1)} .
    \end{equation}
    We now prove $\Fb^p_{n,m} = \Cut^p_{n,m}$ for all $n$, $m$, by induction on $p$.
    For $p = 0$, this is (F4), together with (F4) for $\Cut$.
    Suppose it holds for $p$, and let $M \in K^{(m+p+1) \times (n+p+1)}$.
    Then
    \begin{align*}
        \Fb^{p+1}_{n,m}(M)
         & = \Fb^p_{n,m}\bigl( \Fb^1_{n+p,\,m+p}(M) \bigr)   &  & \text{by (F5) for $\Fb$, with $q = 1$} \\
         & = \Fb^p_{n,m}\bigl( \Cut^1_{n+p,\,m+p}(M) \bigr)  &  & \text{by \cref{eq:one-wire-cut}} \\
         & = \Cut^p_{n,m}\bigl( \Cut^1_{n+p,\,m+p}(M) \bigr) &  & \text{by the induction hypothesis} \\
         & = \Cut^{p+1}_{n,m}(M)                             &  & \text{by (F5) for $\Cut$ (\cref{prop:cut}).}
    \end{align*}
\end{proof}

\begin{remark}[Which hypotheses were used]
    The proof uses (F1), (F2), (F4), (F5), (F6), and it uses (F3) only once: for $p = 2$ and the single matrix $h$ of \cref{lem:functional-equation}.
    It uses that $K$ is a field in \cref{lem:solutions}, where we divide by $\psi(a_0)$.
    The braiding is not used.
\end{remark}

\begin{corollary}[No trace]
    \label{cor:no-trace}
    There is no trace operator on $\Mat_K$ in the sense of \lbl{def:traced-strict-monoidal-cat}.
\end{corollary}

\begin{proof}
    The axioms of a trace operator are the axioms of a feedback operator, with dinaturality required for all morphisms instead of for isomorphisms only, together with the yanking axiom.
    So a trace operator $\Tr$ on $\Mat_K$ would be a feedback operator, and by \cref{thm:unique} it would be $\Cut$.
    The yanking axiom for $K^1$ requires $\Tr^1(S_{1,1}) = I_1$.
    But $\Cut^1(S_{1,1})$ is the block in the upper left of $S_{1,1} = \begin{bsmallmatrix} 0 & 1 \\ 1 & 0 \end{bsmallmatrix}$, which is the matrix with the single entry $0$, and $0 \neq 1$ in a field.
\end{proof}

\begin{corollary}[The delay of the cut]
    \label{cor:cut-delay}
    The delay of the unique feedback operator on $\Mat_K$ is zero:
    \begin{equation}
        \partial_p \coloneqq \Cut^p(S_{p,p}) = 0 \in K^{p \times p} .
    \end{equation}
    For every $M \in K^{(m+p) \times (n+p)}$, the matrices $M \then (\id_{K^m} \oplus \partial_p)$ and $(\id_{K^n} \oplus \partial_p) \then M$ are $p$-traceable, and
    \begin{equation}
        \Cut^p(M) = \Tr^p\bigl( M \then (\id_{K^m} \oplus \partial_p) \bigr) = \Tr^p\bigl( (\id_{K^n} \oplus \partial_p) \then M \bigr) .
    \end{equation}
\end{corollary}

\begin{proof}
    The block in the upper left of \cref{eq:swap} is zero.
    The second statement is the case $\lambda = 0$ of \cref{lem:guard} below, since $\lambda I_p = 0 = \partial_p$ and $\Fb_0^p = \Cut^p$.
\end{proof}

\section{Partial traces guarded by a scalar}
\label{sec:guarded}

The cut is the partial trace guarded by the zero matrix.
We now replace zero by an arbitrary scalar.
The result is no longer defined for all matrices.

\begin{definition}[Guarded partial trace]
    \label{def:guarded}
    Let $\lambda \in K$.
    A matrix $M \in K^{(m+p) \times (n+p)}$, with blocks as in \cref{eq:blocks}, is \emph{$\lambda$-traceable with respect to $p$} if $I_p - \lambda M_{22}$ is invertible.
    In this case we define
    \begin{equation}
        \label{eq:guarded}
        \Fb_\lambda^p(M) \coloneqq M_{11} + \lambda\, M_{12} (I_p - \lambda M_{22})^{-1} M_{21} \in K^{m \times n} .
    \end{equation}
\end{definition}

For $\lambda = 1$ this is \cref{def:tr}: $\Fb_1^p = \Tr^p$.
For $\lambda = 0$, every matrix is $0$-traceable and $\Fb_0^p = \Cut^p$.

\begin{lemma}[$\Fb_\lambda$ is the partial trace guarded by $\lambda$]
    \label{lem:guard}
    Let $\lambda \in K$ and $M \in K^{(m+p) \times (n+p)}$.
    The following are equivalent: $M$ is $\lambda$-traceable; the matrix $M \then (\id_{K^m} \oplus \lambda I_p)$ is $p$-traceable; the matrix $(\id_{K^n} \oplus \lambda I_p) \then M$ is $p$-traceable.
    In this case,
    \begin{equation}
        \label{eq:guard}
        \Fb_\lambda^p(M) = \Tr^p\bigl( M \then (\id_{K^m} \oplus \lambda I_p) \bigr) = \Tr^p\bigl( (\id_{K^n} \oplus \lambda I_p) \then M \bigr) .
    \end{equation}
\end{lemma}

\begin{proof}
    By \cref{eq:mat-operations},
    \begin{align*}
        M \then (\id_{K^m} \oplus \lambda I_p) & = \begin{bmatrix} I_m & 0 \\ 0 & \lambda I_p \end{bmatrix} M = \begin{bmatrix} M_{11} & M_{12} \\ \lambda M_{21} & \lambda M_{22} \end{bmatrix}, \\
        (\id_{K^n} \oplus \lambda I_p) \then M & = M \begin{bmatrix} I_n & 0 \\ 0 & \lambda I_p \end{bmatrix} = \begin{bmatrix} M_{11} & \lambda M_{12} \\ M_{21} & \lambda M_{22} \end{bmatrix} .
    \end{align*}
    Both matrices have the block $\lambda M_{22}$ in the lower right, so each of them is $p$-traceable if and only if $I_p - \lambda M_{22}$ is invertible, that is, if and only if $M$ is $\lambda$-traceable.
    In this case, \cref{eq:tr} gives for the two matrices
    \begin{equation*}
        M_{11} + M_{12} (I_p - \lambda M_{22})^{-1} \lambda M_{21}
        \qquad \text{and} \qquad
        M_{11} + \lambda M_{12} (I_p - \lambda M_{22})^{-1} M_{21},
    \end{equation*}
    and both are equal to \cref{eq:guarded}, because the scalar $\lambda$ commutes with every matrix.
\end{proof}

\begin{lemma}[The delay of $\Fb_\lambda$]
    \label{lem:delay}
    For every $\lambda \in K$ and every $p$, the matrix $S_{p,p}$ is $\lambda$-traceable with respect to $p$, and
    \begin{equation}
        \Fb_\lambda^p(S_{p,p}) = \lambda I_p .
    \end{equation}
\end{lemma}

\begin{proof}
    By \cref{eq:swap}, the blocks of $S_{p,p}$ with respect to $p$ are $M_{11} = 0$, $M_{12} = I_p$, $M_{21} = I_p$, and $M_{22} = 0$.
    So $I_p - \lambda M_{22} = I_p$ is invertible, and \cref{eq:guarded} gives $0 + \lambda\, I_p\, I_p^{-1}\, I_p = \lambda I_p$.
\end{proof}

Together, \cref{lem:guard,lem:delay} say: if we write $\partial^\lambda_p \coloneqq \Fb^p_\lambda(S_{p,p})$ for the delay of $\Fb_\lambda$, then
\begin{equation}
    \label{eq:guarded-headline}
    \Fb^p_\lambda(M) = \Tr^p\bigl( M \then (\id_{K^m} \oplus \partial^\lambda_p) \bigr) = \Tr^p\bigl( (\id_{K^n} \oplus \partial^\lambda_p) \then M \bigr)
\end{equation}
whenever one of the three expressions is defined.
This is the analogue for matrices of the theorem of the first note.
Here it is a reformulation of \cref{def:guarded}; the content is that the guard can be recovered from the operator, as its delay.

\begin{remark}[One wire]
    \label{rem:guarded-one-wire}
    For $p = 1$ and $M$ as in \cref{eq:blocks-one}, the matrix $M$ is $\lambda$-traceable if and only if $\lambda d \neq 1$, and then \cref{eq:guarded} reads
    \begin{equation}
        \label{eq:guarded-one-wire}
        \Fb^1_\lambda(M) = M_{11} + \varphi_\lambda(d)\, c\, r,
        \qquad
        \varphi_\lambda(d) \coloneqq \frac{\lambda}{1 - \lambda d} .
    \end{equation}
    This is the formula of \cref{lem:one-wire}.
    The function $\varphi_\lambda$ is defined on the set $U_\lambda \coloneqq \{ d \in K \mid \lambda d \neq 1 \}$.
    If $\lambda \neq 0$, then $U_\lambda = K \setminus \{ \lambda^{-1} \}$ and $\varphi_\lambda(d) = \frac{1}{\lambda^{-1} - d}$.
    So $\varphi_\lambda$ is the solution of \cref{lem:solutions} with $c = \lambda^{-1}$, and $\varphi_0 = 0$ is the zero solution, with $U_0 = K$.
\end{remark}

\begin{remark}[There are no other candidates, with one exception]
    \label{rem:classification}
    By \cref{lem:solutions}, the functions $\psi$ which satisfy the functional equation on a subset $U \subseteq K$ are the zero function and the functions $a \mapsto \frac{1}{c - a}$ with $c \notin U$.
    By \cref{rem:guarded-one-wire}, for $c \neq 0$ these are restrictions of the functions $\varphi_\lambda$ with $\lambda = c^{-1}$.
    The case $c = 0$ remains: the function $\psi(a) = -a^{-1}$ on $K \setminus \{0\}$.
    Formula \cref{eq:one-wire} with this function reads $M_{11} - d^{-1}\, c\, r$, for $d \neq 0$.
    This is the case $p = 1$ of the operator
    \begin{equation}
        \label{eq:schur}
        \Fb^p_\infty(M) \coloneqq M_{11} - M_{12}\, M_{22}^{-1}\, M_{21},
    \end{equation}
    which is defined if $M_{22}$ is invertible.
    This is the Schur complement of $M_{22}$ in $M$.
    It is not of the form $\Fb_\lambda$ for any $\lambda \in K$; formally it is the case ``$\lambda = \infty$'' of \cref{eq:guarded}.
    The matrix $S_{p,p}$ is not in its domain if $p \geq 1$, so this operator has no delay.

    This suggests that the operators $\Fb_\lambda$, together with $\Fb_\infty$, are the only partially defined operators on $\Mat_K$ which can satisfy the axioms of a feedback operator wherever both sides of an axiom are defined.
    We do not make this statement precise here, and we do not prove it: \cref{lem:one-wire,lem:functional-equation} apply an axiom to matrices such as $K_d$, $L$, and $R$, and for a partially defined operator one has to keep track of which of these are in its domain.
\end{remark}

\section{The transfer function}
\label{sec:transfer}

We use the notation of the first note.
An LTI system $f \colon \R^n \to \R^m$ with state space $\R^k$ is a tuple $f = \sys{\R^k}{A}{B}{C}{D}$ of real matrices $A \in \R^{k \times k}$, $B \in \R^{k \times n}$, $C \in \R^{m \times k}$, and $D \in \R^{m \times n}$.

\begin{definition}[Matrix of an LTI system]
    \label{def:static-matrix}
Let $f = \sys{\R^k}{A}{B}{C}{D}$ be an LTI system from $\R^n$ to $\R^m$.
    The \emph{matrix} of $f$ is
    \begin{equation}
        \label{eq:static-matrix}
        M_f \coloneqq \begin{bmatrix} D & C \\ B & A \end{bmatrix} \in \R^{(m+k) \times (n+k)} .
    \end{equation}
    This is the decomposition of $M_f$ with respect to $k$: its blocks \cref{eq:blocks} are $M_{11} = D$, $M_{12} = C$, $M_{21} = B$, and $M_{22} = A$.
\end{definition}

In terms of the equations $\dot x = A x + B u$, $y = C x + D u$ of the system, $M_f$ is the matrix which maps the input $u$, stacked on top of the state $x$, to the output $y$, stacked on top of $\dot x$.
The wire of the loop is the state.

\begin{proposition}[Evaluating at a point]
    \label{prop:evaluation}
    Let $f = \sys{\R^k}{A}{B}{C}{D} \colon \R^n \to \R^m$ be an LTI system.
    Let $K$ be $\R$ or $\Rs$, and let $c \in K$ with $c \neq 0$.
    Then $M_f$, as a matrix with entries in $K$, is $c^{-1}$-traceable with respect to $k$ if and only if $c I_k - A$ is invertible, and in this case
    \begin{equation}
        \label{eq:evaluation}
        \Fb^k_{1/c}(M_f) = D + C\, (c I_k - A)^{-1} B .
    \end{equation}
\end{proposition}

\begin{proof}
    We have $I_k - c^{-1} A = c^{-1} (c I_k - A)$, and $c^{-1} \neq 0$.
    So $I_k - c^{-1} A$ is invertible if and only if $c I_k - A$ is, and then $(I_k - c^{-1} A)^{-1} = c\, (c I_k - A)^{-1}$.
    By \cref{eq:guarded}, with the blocks of \cref{def:static-matrix},
    \begin{equation*}
        \Fb^k_{1/c}(M_f) = D + c^{-1}\, C \cdot c\, (c I_k - A)^{-1} B = D + C\, (c I_k - A)^{-1} B . \qedhere
    \end{equation*}
\end{proof}

For $K = \R$, the matrix $c I_k - A$ is invertible if and only if $c$ is not an eigenvalue of $A$.
For $K = \Rs$ and $c = s$, the matrix $s I_k - A$ is always invertible: its determinant is a polynomial in $s$ of degree $k$ with leading coefficient $1$, and so it is not the zero element of $\Rs$.

\begin{definition}[Transfer function]
    \label{def:transfer}
Let $f = \sys{\R^k}{A}{B}{C}{D}$ be an LTI system from $\R^n$ to $\R^m$.
    The \emph{transfer function} of $f$ is
    \begin{equation}
        \label{eq:transfer}
        G_f \coloneqq D + C\, (s I_k - A)^{-1} B \in \Rs^{m \times n} .
    \end{equation}
    By \cref{prop:evaluation}, with $K = \Rs$ and $c = s$, the matrix $M_f$ is $s^{-1}$-traceable with respect to $k$, and
    \begin{equation}
        \label{eq:transfer-as-fb}
        G_f = \Fb^k_{1/s}(M_f) .
    \end{equation}
\end{definition}

By \cref{lem:guard}, we can also write \cref{eq:transfer-as-fb} as $G_f = \Tr^k\bigl( M_f \then (\id \oplus s^{-1} I_k) \bigr)$: the transfer function is the partial trace of the matrix of the system, guarded by $s^{-1}$.

\begin{lemma}
    \label{lem:transfer-basic}
    The transfer function has the following two properties.
    \begin{enumerate}[label=\textup{(\alph*)}]
        \item Equivalent LTI systems have the same transfer function: $G_{T \cdot f} = G_f$ for every invertible $T \in \R^{k \times k}$, where $T \cdot f = \sys{\R^k}{T A T^{-1}}{T B}{C T^{-1}}{D}$.
        \item The transfer function of the system $d_p = \sys{\R^p}{0}{I_p}{I_p}{0}$, which represents the delay of $\LTIq$, is $G_{d_p} = s^{-1} I_p$.
    \end{enumerate}
\end{lemma}

\begin{proof}
    (a) We have $s I_k - T A T^{-1} = T (s I_k - A) T^{-1}$, hence $(s I_k - T A T^{-1})^{-1} = T (s I_k - A)^{-1} T^{-1}$, and therefore
    \begin{equation*}
        G_{T \cdot f} = D + C T^{-1} \cdot T (s I_k - A)^{-1} T^{-1} \cdot T B = D + C (s I_k - A)^{-1} B = G_f .
    \end{equation*}

    (b) By \cref{eq:transfer}, $G_{d_p} = 0 + I_p\, (s I_p - 0)^{-1} I_p = s^{-1} I_p$.
\end{proof}

So the transfer function is defined for the morphisms of $\LTIq$, and it maps the delay $\partial_p$ of $\LTIq$ to the delay $s^{-1} I_p$ of $\Fb_{1/s}$ (\cref{lem:delay}).
The main result of this section is that it maps the feedback operator of $\LTIq$ to $\Fb_{1/s}$.
For the proof we use a description of $\Fb_\lambda$ in terms of linear equations.

\begin{lemma}[$\Fb_\lambda$ in terms of equations]
    \label{lem:equations}
    Let $\lambda \in K$ and $M \in K^{(m+p) \times (n+p)}$, with blocks as in \cref{eq:blocks}.
    \begin{enumerate}[label=\textup{(\alph*)}]
        \item The matrix $M$ is $\lambda$-traceable with respect to $p$ if and only if $x = 0$ is the only vector $x \in K^p$ with $x = \lambda M_{22}\, x$.
        \item Suppose that $M$ is $\lambda$-traceable with respect to $p$.
              Let $u \in K^n$, $y \in K^m$, and $x \in K^p$ be vectors with
              \begin{equation}
                  \label{eq:loop-equations}
                  y = M_{11}\, u + M_{12}\, x,
                  \qquad
                  x = \lambda\, ( M_{21}\, u + M_{22}\, x ) .
              \end{equation}
              Then $y = \Fb^p_\lambda(M)\, u$.
    \end{enumerate}
\end{lemma}

\begin{proof}
    (a) The equation $x = \lambda M_{22}\, x$ is equivalent to $(I_p - \lambda M_{22})\, x = 0$, and a square matrix with entries in a field is invertible if and only if the zero vector is the only vector which it maps to zero.

    (b) The second equation of \cref{eq:loop-equations} is equivalent to $(I_p - \lambda M_{22})\, x = \lambda M_{21}\, u$, hence to $x = \lambda (I_p - \lambda M_{22})^{-1} M_{21}\, u$.
    Inserting this into the first equation gives $y = M_{11}\, u + \lambda M_{12} (I_p - \lambda M_{22})^{-1} M_{21}\, u = \Fb^p_\lambda(M)\, u$.
\end{proof}

Recall from the first note the feedback operator of LTI systems.
Let $f \colon \R^{n+p} \to \R^{m+p}$ be an LTI system, with $f = \sys{\R^k}{A}{B}{C}{D}$ and with the blocks
\begin{equation}
    \label{eq:lti-blocks}
    B = \begin{bmatrix} B_1 & B_2 \end{bmatrix},
    \qquad
    C = \begin{bmatrix} C_1 \\ C_2 \end{bmatrix},
    \qquad
    D = \begin{bmatrix} D_{11} & D_{12} \\ D_{21} & D_{22} \end{bmatrix}
\end{equation}
with respect to $p$.
Then $\Fb^p(f)$ is the system
\begin{equation}
    \label{eq:fb-lti}
    \Fb^p(f) = \sys{\R^{p+k}}
    {\begin{bmatrix} D_{22} & C_2 \\ B_2 & A \end{bmatrix}}
    {\begin{bmatrix} D_{21} \\ B_1 \end{bmatrix}}
    {\begin{bmatrix} D_{12} & C_1 \end{bmatrix}}
    {D_{11}}
    \colon \R^n \to \R^m .
\end{equation}

\begin{theorem}[The transfer function maps feedback to guarded partial trace]
    \label{thm:transfer}
    Let $f = \sys{\R^k}{A}{B}{C}{D} \colon \R^{n+p} \to \R^{m+p}$ be an LTI system.
    Then its transfer function $G_f \in \Rs^{(m+p) \times (n+p)}$ is $s^{-1}$-traceable with respect to $p$, and
    \begin{equation}
        \label{eq:transfer-theorem}
        G_{\Fb^p(f)} = \Fb^p_{1/s}(G_f) = \Tr^p\bigl( G_f \then (\id \oplus G_{d_p}) \bigr) .
    \end{equation}
\end{theorem}

\begin{proof}
    We work over the field $K = \Rs$, with $\lambda = s^{-1}$.
    The second equality in \cref{eq:transfer-theorem} holds by \cref{lem:guard} and \cref{lem:transfer-basic}\,(b).
    We prove the first one.

    \emph{Step 1: the systems $f$ and $\Fb^p(f)$ have the same matrix.}
    By \cref{eq:static-matrix} and \cref{eq:lti-blocks},
    \begin{equation}
        \label{eq:N}
        M_f = \begin{bmatrix} D & C \\ B & A \end{bmatrix}
        = \begin{bmatrix} D_{11} & D_{12} & C_1 \\ D_{21} & D_{22} & C_2 \\ B_1 & B_2 & A \end{bmatrix}
        \eqqcolon N,
    \end{equation}
    where the rows are grouped as $m + p + k$ and the columns as $n + p + k$.
    By \cref{eq:static-matrix} and \cref{eq:fb-lti}, the matrix of $\Fb^p(f)$ is
    \begin{equation*}
        M_{\Fb^p(f)} =
        \left[\begin{array}{c|cc}
                D_{11} & D_{12} & C_1 \\ \hline
                D_{21} & D_{22} & C_2 \\
                B_1    & B_2    & A
            \end{array}\right]
        = N .
    \end{equation*}
    The two systems differ only in which part of $N$ is regarded as the loop: for $f$ it is the last $k$ rows and columns, for $\Fb^p(f)$ it is the last $p + k$ rows and columns.
    By \cref{eq:transfer-as-fb}, applied to $f$ and to $\Fb^p(f)$, whose state space is $\R^{p+k}$,
    \begin{equation}
        \label{eq:two-transfer-functions}
        G_f = \Fb^k_{1/s}(N),
        \qquad
        G_{\Fb^p(f)} = \Fb^{p+k}_{1/s}(N),
    \end{equation}
    and in particular $N$ is $s^{-1}$-traceable with respect to $p + k$.
    So we have to prove that $G_f$ is $s^{-1}$-traceable with respect to $p$, and that $\Fb^p_{1/s}\bigl( \Fb^k_{1/s}(N) \bigr) = \Fb^{p+k}_{1/s}(N)$.

    \emph{Step 2: notation.}
    Set $X \coloneqq (s I_k - A)^{-1} \in \Rs^{k \times k}$.
    By \cref{eq:transfer} and \cref{eq:lti-blocks}, the blocks of $G \coloneqq G_f$ with respect to $p$ are
    \begin{equation}
        \label{eq:G-blocks}
        G_{11} = D_{11} + C_1 X B_1,
        \quad
        G_{12} = D_{12} + C_1 X B_2,
        \quad
        G_{21} = D_{21} + C_2 X B_1,
        \quad
        G_{22} = D_{22} + C_2 X B_2 .
    \end{equation}
    The blocks of $N$ with respect to $p + k$ are
    \begin{equation}
        \label{eq:N-blocks}
        N_{11} = D_{11},
        \qquad
        N_{12} = \begin{bmatrix} D_{12} & C_1 \end{bmatrix},
        \qquad
        N_{21} = \begin{bmatrix} D_{21} \\ B_1 \end{bmatrix},
        \qquad
        N_{22} = \begin{bmatrix} D_{22} & C_2 \\ B_2 & A \end{bmatrix} .
    \end{equation}
    We also note: for vectors $v, x \in K^k$,
    \begin{equation}
        \label{eq:X}
        x = X v
        \quad \Longleftrightarrow \quad
        (s I_k - A)\, x = v
        \quad \Longleftrightarrow \quad
        x = s^{-1} ( v + A x ) .
    \end{equation}

    \emph{Step 3: $G$ is $s^{-1}$-traceable with respect to $p$.}
    By \cref{lem:equations}\,(a), we have to show: if $w \in K^p$ satisfies $w = s^{-1} G_{22}\, w$, then $w = 0$.
    Let $w$ be such a vector, and set $x \coloneqq X B_2\, w$.
    By \cref{eq:X}, with $v = B_2\, w$,
    \begin{equation*}
        x = s^{-1} ( B_2\, w + A x ) .
    \end{equation*}
    By \cref{eq:G-blocks}, $G_{22}\, w = D_{22}\, w + C_2 X B_2\, w = D_{22}\, w + C_2\, x$, and so
    \begin{equation*}
        w = s^{-1} G_{22}\, w = s^{-1} ( D_{22}\, w + C_2\, x ) .
    \end{equation*}
    The last two displayed equations say, by \cref{eq:N-blocks},
    \begin{equation*}
        \begin{bmatrix} w \\ x \end{bmatrix} = s^{-1} N_{22} \begin{bmatrix} w \\ x \end{bmatrix} .
    \end{equation*}
    Since $N$ is $s^{-1}$-traceable with respect to $p + k$, \cref{lem:equations}\,(a) gives $\begin{bsmallmatrix} w \\ x \end{bsmallmatrix} = 0$, and in particular $w = 0$.

    \emph{Step 4: the two matrices agree.}
    Let $u \in K^n$ be an arbitrary vector, and set
    \begin{gather*}
        w \coloneqq s^{-1} (I_p - s^{-1} G_{22})^{-1} G_{21}\, u \in K^p, \\
        y \coloneqq G_{11}\, u + G_{12}\, w \in K^m,
        \qquad
        x \coloneqq X ( B_1\, u + B_2\, w ) \in K^k .
    \end{gather*}
    The inverse exists by Step 3.
    By the definition of $w$, and by \cref{eq:guarded} for $G$,
    \begin{equation}
        \label{eq:w-and-y}
        w = s^{-1} ( G_{21}\, u + G_{22}\, w ),
        \qquad
        y = \Fb^p_{1/s}(G)\, u .
    \end{equation}
    By \cref{eq:X}, with $v = B_1\, u + B_2\, w$,
    \begin{equation}
        \label{eq:x-equation}
        x = s^{-1} ( B_1\, u + B_2\, w + A x ) .
    \end{equation}
    By \cref{eq:G-blocks} and the definition of $x$,
    \begin{align*}
        G_{11}\, u + G_{12}\, w & = D_{11}\, u + D_{12}\, w + C_1 X ( B_1\, u + B_2\, w ) = D_{11}\, u + D_{12}\, w + C_1\, x, \\
        G_{21}\, u + G_{22}\, w & = D_{21}\, u + D_{22}\, w + C_2 X ( B_1\, u + B_2\, w ) = D_{21}\, u + D_{22}\, w + C_2\, x .
    \end{align*}
    Inserting this into the definition of $y$ and into the first equation of \cref{eq:w-and-y}, we obtain
    \begin{equation}
        \label{eq:y-and-w}
        y = D_{11}\, u + D_{12}\, w + C_1\, x,
        \qquad
        w = s^{-1} ( D_{21}\, u + D_{22}\, w + C_2\, x ) .
    \end{equation}
    By \cref{eq:N-blocks}, the equations \cref{eq:y-and-w} and \cref{eq:x-equation} together say
    \begin{equation*}
        y = N_{11}\, u + N_{12} \begin{bmatrix} w \\ x \end{bmatrix},
        \qquad
        \begin{bmatrix} w \\ x \end{bmatrix} = s^{-1} \Bigl( N_{21}\, u + N_{22} \begin{bmatrix} w \\ x \end{bmatrix} \Bigr) .
    \end{equation*}
    These are the equations \cref{eq:loop-equations} for the matrix $N$, decomposed with respect to $p + k$.
    Since $N$ is $s^{-1}$-traceable with respect to $p + k$, \cref{lem:equations}\,(b) gives $y = \Fb^{p+k}_{1/s}(N)\, u$.
    Together with the second equation of \cref{eq:w-and-y}, this shows
    \begin{equation*}
        \Fb^p_{1/s}(G)\, u = \Fb^{p+k}_{1/s}(N)\, u \qquad \text{for all } u \in K^n,
    \end{equation*}
    and so $\Fb^p_{1/s}(G_f) = \Fb^{p+k}_{1/s}(N) = G_{\Fb^p(f)}$, by \cref{eq:two-transfer-functions}.
\end{proof}

\begin{remark}[The operators on real matrices are evaluations of the transfer function]
    \label{rem:table}
    By \cref{prop:evaluation}, the partially defined operators of \cref{sec:guarded}, applied to the real matrix $M_f$, are the values of the transfer function at real numbers.
    \begin{center}
        \begin{tabular}{llll}
            \toprule
            $c$           & operator on $\Mat_\R$                 & value on $M_f$              & defined if                        \\
            \midrule
            ``$\infty$''  & $\Cut^k = \Fb^k_0$                    & $D$                         & always                            \\
            $c \neq 0$    & $\Fb^k_{1/c}$                         & $D + C (c I_k - A)^{-1} B$  & $c$ is not an eigenvalue of $A$   \\
            $1$           & $\Tr^k = \Fb^k_1$                     & $D + C (I_k - A)^{-1} B$    & $1$ is not an eigenvalue of $A$   \\
            $0$           & $\Fb^k_\infty$ of \cref{eq:schur}     & $D - C A^{-1} B$            & $A$ is invertible                 \\
            \bottomrule
        \end{tabular}
    \end{center}
    Two of these values have a direct meaning for the system.
    In the discrete-time reading $x_{i+1} = A x_i + B u_i$, $y_i = C x_i + D u_i$, suppose that the input and the state are constant, $u_i = u$ and $x_i = x$.
    Then $x = A x + B u$, and if $1$ is not an eigenvalue of $A$, the output is the constant $y = \bigl( D + C (I_k - A)^{-1} B \bigr) u$.
    So the partial trace of \lbl{def:trace-for-mat}, applied to $M_f$, describes the system in equilibrium.
    This is the first of the two derivations of the formula for the partial trace in the remark after \lbl{def:trace-for-mat}.
    In the continuous-time reading, a constant input $u$ and a constant state $x$ satisfy $0 = A x + B u$, and if $A$ is invertible, the output is the constant $y = ( D - C A^{-1} B ) u$.

    Every real number $c$ is an eigenvalue of some real matrix, so none of the operators $\Fb_{1/c}$ with $c \neq 0$ is defined for all real matrices, in accordance with \cref{thm:unique}.
    The indeterminate $s$ is not an eigenvalue of any real matrix, and this is why $\Fb_{1/s}$ is defined on $M_f$ for every $f$.
\end{remark}

\begin{remark}[What is not proved here]
    \label{rem:not-proved}
    The following statements are not proved in this note.
    \begin{enumerate}[label=\textup{(\arabic*)}]
        \item That $\Fb_\lambda$ satisfies the axioms (F1) to (F6) wherever both sides are defined, and the classification suggested in \cref{rem:classification}.
        \item That $\Fb_{1/s}$ is a feedback operator on matrices of proper rational functions.
              In more detail: \cref{thm:unique} applies to the field $K = \Rs$ as well: the only feedback operator on $\Mat_{\Rs}$ is the cut.
              Accordingly, $\Fb_{1/s}$ is not defined for all matrices with entries in $\Rs$; for example, the matrix $K_s$ of \cref{eq:Kd} is not $s^{-1}$-traceable, since $1 - s^{-1} s = 0$.
              \Cref{thm:transfer} shows that $\Fb_{1/s}$ is defined on all transfer functions.
              The entries of a transfer function are proper rational functions, that is, quotients of polynomials in which the degree of the numerator is at most the degree of the denominator.
              The proper rational functions form a ring which is not a field: $s^{-1}$ is proper, but its inverse $s$ is not.
              We expect that $\Fb_{1/s}$ is defined for all matrices with entries in this ring, and that it is a feedback operator on the category of such matrices.
              This would not contradict \cref{thm:unique}, whose proof divides by $\psi(a_0)$ in \cref{lem:solutions}: for $\psi(a) = \frac{1}{s - a}$, the element $\psi(a)^{-1} = s - a$ is not proper.
              We have not verified this expectation.
        \item That the transfer function is a functor from $\LTIq$ to matrices of proper rational functions which preserves the monoidal product.
              If this and the expectation in (2) are correct, then \cref{thm:transfer} and \cref{lem:transfer-basic} say that it preserves feedback operators and delays.
              We also do not discuss in which sense it is unique with these properties.
    \end{enumerate}
\end{remark}

\end{document}
